\documentclass[12pt,leqno]{article}

\usepackage{amsmath,amssymb,amsthm,colortbl,xcolor,nicefrac,setspace}
\usepackage[T1]{fontenc}
\usepackage[expert,seriftt,lucidasmallscale]{lucidabr}
\usepackage{mathrsfs,hyperref,color,graphicx}
\usepackage{booktabs,tabularx,array,needspace}
\hypersetup{hidelinks,hypertexnames=false,
  pdftitle={Recombination without Exhaustion},
  pdfsubject={Logical Purity, Application and Purity of Pure refute a universal proposition}}

% Self-contained Forbes / Phil-211 / Lucida setup.
\setlength{\textwidth}{6.5in}
\setlength{\textheight}{9in}
\setlength{\oddsidemargin}{0mm}
\setlength{\evensidemargin}{0mm}
\setlength{\topmargin}{0mm}
\setlength{\headheight}{0mm}
\setlength{\headsep}{0mm}
\setlength{\topskip}{0mm}
\setlength{\hoffset}{0in}
\setlength{\voffset}{0in}

\newcommand{\AND}{\ensuremath{\mathbin{\&}}}
\newcommand{\OR}{\ensuremath{\vee}}
\newcommand{\IFF}{\ensuremath{\leftrightarrow}}
\newcommand{\IC}{\ensuremath{\rightarrow}}
\newcommand{\NOT}{\ensuremath{\text{$\thicksim$}}}
\newcommand{\dentails}{\ensuremath{\vdash}}
\newcommand{\Pure}{\ensuremath{\mathrm{Pure}}}
\newcommand{\Fun}{\ensuremath{\mathrm{Fun}}}
\newcommand{\Uni}{\ensuremath{\mathrm{Uni}}}
\newcommand{\Cl}{\ensuremath{\mathrm{Cl}}}
\newcommand{\Pres}{\ensuremath{\mathrm{Pres}}}
\newcommand{\PPu}{\ensuremath{\mathrm{PP}_u}}
\newcommand{\HPE}{\ensuremath{\mathrm{H}+\mathrm{PE}}}
\newcommand{\App}{\ensuremath{\mathrm{App}}}
\newcommand{\Rec}{\ensuremath{\mathrm{Rec}_1}}
\newcommand{\hW}{\ensuremath{\widehat{W}}}
\newcommand{\hCl}{\ensuremath{\widehat{\Cl}}}
\newcommand{\hR}{\ensuremath{\widehat{R}}}
\newcommand{\hF}{\ensuremath{\widehat{F}}}
\newcommand{\hQ}{\ensuremath{\widehat{Q}}}
\newcommand{\hP}{\ensuremath{\widehat{P}}}
\newcommand{\hPp}{\ensuremath{\widehat{P}'}}
\newcommand{\hPhi}{\ensuremath{\widehat{\Phi}}}
\newcommand{\conj}{\ensuremath{\mathsf{and}}}
\newcommand{\cmp}{\ensuremath{\mathsf{comp}}}
\newcommand{\negop}{\ensuremath{\mathsf{neg}}}

\theoremstyle{plain}
\newtheorem{theorem}{Theorem}
\newtheorem{lemma}{Lemma}
\newtheorem*{corollary}{Corollary}
\newtheorem*{proposition}{Proposition}
\renewcommand{\qedsymbol}{\(\blacksquare\)}
\onehalfspacing

\begin{document}

\begin{center}
{\Large Recombination without Exhaustion}

\smallskip
Purity of Pure and fundamental propositions\\
11 November 2026
\end{center}

\noindent
Can purity itself be pure if there is a fundamental proposition? Under
Bacon's full Quantified Logical Necessity, Lettie's theorem answers
Goodman's question negatively. His proof uses two instances of zeroary
Exhaustion. Shortly after Lettie found his result, a strengthening was
found that drops Exhaustion. This note proves it.

The strengthening says that Logical Purity, Application and Purity of
Pure at type \(t\IC t\), taken as ordinary premises over \HPE, already
rule out every universal proposition. Since unary Recombination makes a
fundamental proposition universal, Purity of Pure is refuted by Logical
Purity, Application, unary Recombination and a fundamental proposition.
This answers the Recombination-only reading of Goodman's question. No
Exhaustion, Persistence or Purity of Fun is used.

\section{From fundamentality to a universal proposition}

Propositions have type \(t\), \(u:=t\IC t\) is the type of unary
propositional operators, and \(c:=u\IC t\). Write \(\pi:=\Pure_u\), a
constant of type \(c\). The instance of Purity of Pure at issue is
\[
 \PPu:=\Pure_c(\pi).
\]
Throughout, \(p,q,x,y,z,r,b\) range over propositions and \(S,X,Y\) over
operators of type \(u\). We use the definitions
\[
 \top:=\forall p(p\IC p),\qquad \Box A:=(A=\top),\qquad
 \Diamond A:=\NOT\Box\NOT A .
\]
Following Lettie, put
\begin{align}
 R(x,y)&:=\forall Y\bigl(\pi(Y)\IC(Y(y)\IC\Diamond Y(x))\bigr),
   \label{eq:R}\\
 \Uni(x)&:=\forall y\,R(x,y). \label{eq:Uni}
\end{align}
The relation \(R(x,y)\) says that \(x\) can have every pure property that
\(y\) has, where a property is an operator of type \(u\). We call \(x\)
\emph{universal} if this holds for every proposition \(y\).

Unary Recombination says
\[
 \Rec:\quad\forall X\forall r\bigl((\Pure_u(X)\AND\Fun_t(r))
   \IC(\Box X(r)\IC\forall z\,X(z))\bigr).
\]
Logical Purity says that closed terms containing no constants are pure.
Application at \((\sigma,\tau)\) says that a pure operator applied to a
pure argument has a pure value:
\[
 \App_{\sigma,\tau}:\quad\forall X^{\sigma\IC\tau}\forall y^{\sigma}
   \bigl((\Pure(X)\AND\Pure(y))\IC\Pure(X(y))\bigr).
\]

\begin{proposition}[Lettie: fundamental propositions are universal]
Put \(\negop:=\lambda Y\lambda x.\NOT Y(x)\). Given \(\Pure(\negop)\),
\(\App_{u,u}\) and \(\Rec\), every fundamental proposition is universal.
So if a fundamental proposition exists, a universal one exists.
\end{proposition}
\begin{proof}
Let \(r\) be fundamental. Fix a proposition \(y\) and a pure operator
\(Y\) such that \(Y(y)\). The operator \(X:=\negop\,Y\) is pure, by
\(\Pure(\negop)\), \(\Pure_u(Y)\) and \(\App_{u,u}\). By
\(\beta\)-conversion, \(\NOT Y(r)\IC X(r)\) is a theorem, so
Necessitation and K give \(\Box\NOT Y(r)\IC\Box X(r)\).

If \(\Box\NOT Y(r)\), Recombination at \(X\) gives \(\forall z\,X(z)\),
hence \(X(y)\), that is, \(\NOT Y(y)\), contradicting \(Y(y)\). Hence
\(\Diamond Y(r)\). Since \(y\) and \(Y\) were arbitrary, \(\Uni(r)\).
The second claim follows by existential instantiation on the
fundamental proposition and existential generalization.
\end{proof}

It therefore suffices to show that \(\PPu\), together with suitable
instances of Logical Purity and Application, rules out every universal
proposition. The proof of Theorem~\ref{thm:main} below contains no
\(\Fun\).

\section{The base logic and the idea of the proof}

We work over \HPE: classical higher-order logic H, with Leibniz identity
and \(\beta\eta\)-conversion, closed under modus ponens, Generalization,
existential Instantiation, and Propositional Equivalence at type \(t\).
The last rule takes a \emph{theorem} \(A\IFF B\) to \(A=B\). Modalized
Functionality is not assumed. We write C for Bacon's Classicism, which
extends \HPE.

The base logic proves the laws
\[
 \begin{aligned}
 \mathrm{K}:&\quad \Box(A\IC B)\IC(\Box A\IC\Box B),\\
 \mathrm{T}:&\quad \Box A\IC A,
 &\mathrm{4}:&\quad \Box A\IC\Box\Box A,
 \end{aligned}
\]
and \(\Box(A\IFF B)\IC A=B\). It is closed under Necessitation; hence,
by K, a theorem \(A\IC B\) yields the theorem \(\Box A\IC\Box B\). We call this the K-rule. These rules
apply to base theorems, including theorems that mention the constants
\(\Pure_\sigma\). They do \emph{not} authorize necessitating
consequences of the purity premises: every premise is ordinary, and
``under premises \(\Pi\), \(A\)'' means that \(\bigwedge\Pi\IC A\) is a
theorem of \HPE. We mark the points where premises enter.

\paragraph{Where Lettie uses Exhaustion.}
Lettie calls \(z\) \emph{firm} when
\(F(z):=\Box\bigl(\NOT R(z,\NOT z)\AND\NOT R(\NOT z,z)\bigr)\). He sets
\(Q(z):=F(z)\IC z\) and introduces two tests,
\(P(z):=\NOT z\AND\NOT F(z)\) and \(P'(z):=z\AND\NOT F(z)\). The heart
of his proof is that \(\Box\pi(P)\) and \(\Box\pi(P')\) make every value
of \(Q\) firm. The purity premises supply \(\pi(P)\) and \(\pi(P')\),
but only as consequences of premises, so they cannot be necessitated.
Zeroary Exhaustion, \(\Pure_t(a)\IC(a\IC\Box a)\), supplies the boxes.

\paragraph{The plan.}
We relativize Lettie's construction to a \emph{guard} proposition \(b\):
in \(R\), the possibility \(\Diamond Y(x)\) becomes
\(\Diamond(b\AND Y(x))\), and firmness becomes
\(\Box(b\IC\cdots)\). The relativized argument needs, in place of
\(\Box\pi(P)\AND\Box\pi(P')\), only
\[
 \Box\bigl(b\IC\kappa(b)\bigr),
\]
where \(\kappa(b)\) says that the two relativized tests are pure.
This can be a theorem of the base logic even though \(\kappa(b)\) itself
is only a consequence of premises. We therefore look for a proposition
\(a\) such that
\begin{equation}
 \text{\(a\) is pure and true, under the premises,}
 \qquad\text{and}\qquad
 \dentails\Box\bigl(a\IC\kappa(a)\bigr).
 \label{eq:goals}
\end{equation}
Section~\ref{sec:closure} constructs, for any closed operator \(\Theta\),
a proposition \(a\) such that \(a\IC\Theta(a)\) is a base theorem.
Section~\ref{sec:purity} shows that the premises make \(a\) pure and
true whenever \(\Theta\) is pure and preserves pure truths. Taking \(\Theta(b):=b\AND\kappa(b)\)
gives \eqref{eq:goals}. Lettie's diagonal then runs relative to \(a\):
if \(r\) were universal, \(e:=Q_a(r)\) would satisfy both a separation
condition and the relation it forbids.

\section{A closure lemma for pure truths}\label{sec:closure}

The construction in this section is independent of Recombination and of
Lettie's construction. Fix a closed term \(\Theta\) of type \(u\); it
may contain constants. Let \(m:=\lambda S.\forall p(S(p)\IC p)\), so that
\(m(S)\) says that everything \(S\) holds of is true. For \(X\)
of type \(u\), let \(\Cl(X)\) be the conjunction of the following five
conditions:
\begin{align}
 &X(\top); \tag{Cl1}\\
 &\forall p\forall q\bigl((X(p)\AND X(q))\IC X(p\AND q)\bigr); \tag{Cl2}\\
 &\forall p\bigl(X(p)\IC X(\Theta(p))\bigr); \tag{Cl3}\\
 &\forall S\bigl(\pi(S)\IC X(\pi(S))\bigr); \tag{Cl4}\\
 &\forall S\bigl((\pi(S)\AND\forall p(S(p)\IC X(p)))\IC X(m(S))\bigr).
   \tag{Cl5}
\end{align}
Put
\[
 W:=\lambda q.\forall X\bigl(\Cl(X)\IC X(q)\bigr),\qquad
 a:=m(W)\AND\pi(W).
\]
Thus \(W\) holds of the propositions that belong to every operator
satisfying (Cl1)--(Cl5), and \(a\) says that \(W\) holds only of truths
and is pure. With \(\Theta:=\Phi\), \(a\) will be the guard of
\eqref{eq:goals}. Two facts about \(W\) are used throughout. By
\emph{elimination}, \(W(q)\IC(\Cl(X)\IC X(q))\). By \emph{introduction},
from \(P\IC(\Cl(X)\IC X(q))\), with \(X\) not free in \(P\), we may infer
\(P\IC W(q)\).

The terms \(\Cl\) and \(W\) mention \(\pi\) and \(\Theta\). Abstracting
both gives closed terms that contain no constants:
\[
 \hCl:=\lambda c'\lambda T\lambda X.\,\bigl(\text{(Cl1)}\AND\cdots\AND
   \text{(Cl5)}\bigr),\qquad
 \hW:=\lambda c'\lambda T\lambda q.\forall X\bigl(\hCl\,c'\,T\,X\IC X(q)\bigr),
\]
where \(c'\) and \(T\) stand for \(\pi\) and \(\Theta\) in the
conditions. By \(\beta\)-conversion, \(\Cl(X)=\hCl\,\pi\,\Theta\,X\) and
\(W=\hW\,\pi\,\Theta\). Likewise \(m\) and the conjunction operator
\(\conj:=\lambda p\lambda q.(p\AND q)\) contain no constants.

\begin{lemma}[\(W\) satisfies the closure conditions]\label{lem:ClW}
\(\dentails\Cl(W)\).
\end{lemma}
\begin{proof}
Each condition is proved for \(W\) by introduction, from the same
condition for an arbitrary \(X\) such that \(\Cl(X)\). For (Cl1),
\(\Cl(X)\) gives \(X(\top)\), so \(W(\top)\). For (Cl2), suppose
\(W(p)\AND W(q)\) and \(\Cl(X)\); elimination gives \(X(p)\) and
\(X(q)\), and (Cl2) for \(X\) gives \(X(p\AND q)\); now generalize \(p\)
and \(q\). For (Cl3), elimination turns \(W(p)\) into \(X(p)\), and
(Cl3) for \(X\) gives \(X(\Theta(p))\). For (Cl4), (Cl4) for \(X\) turns
\(\pi(S)\) into \(X(\pi(S))\). For (Cl5), suppose
\(\pi(S)\AND\forall p(S(p)\IC W(p))\) and \(\Cl(X)\). For each \(p\),
\(S(p)\IC W(p)\) and elimination give \(S(p)\IC X(p)\). Generalizing,
\(\forall p(S(p)\IC X(p))\), and (Cl5) for \(X\) gives \(X(m(S))\).
\end{proof}

\begin{lemma}[\(a\IC\Theta(a)\) is a base theorem]\label{lem:base}
\(\dentails a\IC\Theta(a)\), and hence
\(\dentails\Box\bigl(a\IC\Theta(a)\bigr)\).
\end{lemma}
\begin{proof}
Instantiate the conditions of Lemma~\ref{lem:ClW} at \(W\) itself.
Condition (Cl4) at \(S:=W\) gives \(\pi(W)\IC W(\pi(W))\). Condition
(Cl5) at \(S:=W\), with the theorem \(\forall p(W(p)\IC W(p))\), gives
\(\pi(W)\IC W(m(W))\). Condition (Cl2) gives
\((W(m(W))\AND W(\pi(W)))\IC W(a)\), and (Cl3) gives
\(W(a)\IC W(\Theta(a))\). Finally, \(m(W)\IFF\forall p(W(p)\IC p)\) by
\(\beta\)-conversion, so \(m(W)\IC(W(\Theta(a))\IC\Theta(a))\).

Now assume \(a\), that is, \(m(W)\) and \(\pi(W)\). The first three
facts give \(W(a)\), the fourth gives \(W(\Theta(a))\), and the last
gives \(\Theta(a)\). No premise was used, so \(a\IC\Theta(a)\) is a
theorem of \HPE, and Necessitation applies.
\end{proof}

The key condition is (Cl4). At \(S=W\) it places the purity claim \(\pi(W)\)
inside \(W\), and this is what lets \(a\IC\Theta(a)\) hold in the base
logic. Condition (Cl1) is not needed here.

\section{The relativized Lettie construction}\label{sec:relative}

For a proposition \(b\), define
\begin{align}
 R_b(x,y)&:=\forall Y\bigl(\pi(Y)\IC(Y(y)\IC\Diamond(b\AND Y(x)))\bigr),
   \label{eq:Rb}\\
 C_b(z)&:=\NOT R_b(z,\NOT z)\AND\NOT R_b(\NOT z,z), \label{eq:Cb}\\
 F_b(z)&:=\Box\bigl(b\IC C_b(z)\bigr), \label{eq:Fb}\\
 Q_b(z)&:=F_b(z)\IC z, \label{eq:Qb}\\
 P_b(z)&:=\NOT z\AND\NOT F_b(z),\qquad P'_b(z):=z\AND\NOT F_b(z),
   \label{eq:Pb}\\
 \kappa(b)&:=\pi(P_b)\AND\pi(P'_b),\qquad
 \Phi(b):=b\AND\kappa(b). \label{eq:kappa}
\end{align}
The relation \(R_b(x,y)\) says that \(x\) can have, compatibly with
\(b\), every pure property that \(y\) has. Call \(z\)
\emph{\(b\)-firm} when \(F_b(z)\): necessarily, if \(b\), then neither
\(z\) nor \(\NOT z\) can have every pure property of the other
compatibly with \(b\). Up to Boolean identities, the case \(b=\top\) is
Lettie's construction.

Each of these expressions is built from logical vocabulary, \(\pi\) and
\(b\). Abstracting \(\pi\) and \(b\) gives closed terms that contain no
constants:
\begin{align*}
 \hR&:=\lambda c'\lambda b\lambda x\lambda y.\,
   \forall Y\bigl(c'(Y)\IC(Y(y)\IC\Diamond(b\AND Y(x)))\bigr),\\
 \hF&:=\lambda c'\lambda b\lambda z.\,
   \Box\bigl(b\IC(\NOT\hR c'bz(\NOT z)\AND\NOT\hR c'b(\NOT z)z)\bigr),\\
 \hQ&:=\lambda c'\lambda b\lambda z.(\hF c'bz\IC z),\\
 \hP&:=\lambda c'\lambda b\lambda z.(\NOT z\AND\NOT\hF c'bz),\qquad
 \hPp:=\lambda c'\lambda b\lambda z.(z\AND\NOT\hF c'bz),\\
 \hPhi&:=\lambda c'\lambda b.\bigl(b\AND(c'(\hP c'b)\AND c'(\hPp c'b))\bigr).
\end{align*}
By \(\beta\)-conversion, \(R_b(x,y)=\hR\pi bxy\), \(F_b(z)=\hF\pi bz\),
\(Q_b=\hQ\pi b\), \(P_b=\hP\pi b\), \(P'_b=\hPp\pi b\) and
\(\Phi=\hPhi\pi\). For \(R_b\) we use elimination, its instance at a
given \(Y\), and introduction: from
\(P\IC(\pi(Y)\IC(Y(y)\IC\Diamond(b\AND Y(x))))\), with \(Y\) not free
in \(P\), infer \(P\IC R_b(x,y)\).

All results in this section are theorems of \HPE, for arbitrary
propositions \(b\), \(x\) and \(z\). They mention \(\pi\) and \(b\) but
make no assumption about what is pure or true.

\begin{lemma}[Firmness persists]\label{lem:A}
\(F_b(x)\IC\Box F_b(x)\).
\end{lemma}
\begin{proof}
By definition \(F_b(x)\) is \(\Box B\) with \(B:=b\IC C_b(x)\). Axiom 4
gives \(\Box B\IC\Box\Box B\), and the K-rule applied to
\(\Box B\IC F_b(x)\) gives \(\Box\Box B\IC\Box F_b(x)\).
\end{proof}

\begin{lemma}[\(Q_b\) fixes firm propositions]\label{lem:B}
\(F_b(x)\IC Q_b(x)=x\).
\end{lemma}
\begin{proof}
The formula \(F_b(x)\IC(Q_b(x)\IFF x)\) is a tautology, given
\eqref{eq:Qb}. The K-rule gives
\(\Box F_b(x)\IC\Box(Q_b(x)\IFF x)\). Assume \(F_b(x)\).
Lemma~\ref{lem:A} supplies \(\Box F_b(x)\), so
\(\Box(Q_b(x)\IFF x)\), and the law \(\Box(p\IFF q)\IC p=q\) gives
\(Q_b(x)=x\).
\end{proof}

\begin{lemma}[Symmetry]\label{lem:C}
\(F_b(\NOT x)\IFF F_b(x)\), and \(P'_b(z)=P_b(\NOT z)\).
\end{lemma}
\begin{proof}
Propositional Equivalence gives \(\NOT\NOT x=x\). By congruence,
\(R_b(\NOT x,\NOT\NOT x)\IFF R_b(\NOT x,x)\) and
\(R_b(\NOT\NOT x,\NOT x)\IFF R_b(x,\NOT x)\). Hence
\(C_b(\NOT x)\IFF C_b(x)\), with the two conjuncts exchanged, and
Propositional Equivalence gives \(C_b(\NOT x)=C_b(x)\). Congruence in
\(\Box(b\IC\cdot)\) gives the first claim.

For the second, \(P'_b(z)\IFF(z\AND\NOT F_b(z))\) and
\(P_b(\NOT z)\IFF(\NOT\NOT z\AND\NOT F_b(\NOT z))\). By the first claim
the right-hand sides are equivalent, and Propositional Equivalence gives
the identity.
\end{proof}

The choice \(Q_b(x)=F_b(x)\IC x\) has a useful consequence: a false
value of \(Q_b\) is \(b\)-firm.

\begin{lemma}\label{lem:D}
\(\NOT Q_b(x)\IC F_b(Q_b(x))\).
\end{lemma}
\begin{proof}
Assume \(\NOT Q_b(x)\), that is, \(\NOT(F_b(x)\IC x)\). Then \(F_b(x)\),
and Lemma~\ref{lem:B} gives \(Q_b(x)=x\). Substituting \(Q_b(x)\) for
\(x\) in \(F_b(x)\) by Leibniz's law yields \(F_b(Q_b(x))\).
\end{proof}

\begin{lemma}[The first test necessarily fails on \(Q_b\)'s values]\label{lem:E}
\(\Box\NOT P_b(Q_b(x))\).
\end{lemma}
\begin{proof}
Put \(q:=Q_b(x)\). Then \(P_b(q)\IFF(\NOT q\AND\NOT F_b(q))\), and
Lemma~\ref{lem:D} gives \(\NOT q\IC F_b(q)\). Thus \(\NOT P_b(q)\) is a
base theorem, and Necessitation gives the result.
\end{proof}

The next lemma is the relativized form of Lettie's key step. The guard
appears in it unboxed.

\begin{lemma}[The guarded step]\label{lem:guarded}
\(\bigl(b\AND\kappa(b)\bigr)\IC C_b(Q_b(x))\).
\end{lemma}
\begin{proof}
Put \(q:=Q_b(x)\). Assume \(b\), \(\pi(P_b)\) and \(\pi(P'_b)\), and
suppose, for a contradiction, that \(C_b(q)\) is false. Thus
\[
 R_b(q,\NOT q)\OR R_b(\NOT q,q).
\]
The proposition \(q\) is not \(b\)-firm: \(F_b(q)\) would give
\(b\IC C_b(q)\) by T, and so \(C_b(q)\). Since \(\NOT F_b(q)\), the
contrapositive of Lemma~\ref{lem:D} gives \(q\). Consequently
\[
 P'_b(q),\qquad P_b(\NOT q).
\]
The first follows from \(q\AND\NOT F_b(q)\). For the second,
\(P_b(\NOT q)\IFF(\NOT\NOT q\AND\NOT F_b(\NOT q))\); here \(q\) gives
\(\NOT\NOT q\), and \(F_b(\NOT q)\IFF F_b(q)\) by Lemma~\ref{lem:C}.

If \(R_b(q,\NOT q)\), instantiate its quantifier at the pure operator
\(P_b\). Since \(P_b(\NOT q)\), it gives \(\Diamond(b\AND P_b(q))\).
If instead \(R_b(\NOT q,q)\), instantiate at the pure operator
\(P'_b\). Since \(P'_b(q)\), it gives \(\Diamond(b\AND P'_b(\NOT q))\).
But Lemma~\ref{lem:C} at \(\NOT q\) gives
\(P'_b(\NOT q)=P_b(\NOT\NOT q)\), and \(\NOT\NOT q=q\) gives
\(P_b(\NOT\NOT q)\IFF P_b(q)\). Hence \(P'_b(\NOT q)\IFF P_b(q)\);
Propositional Equivalence makes this an identity, and congruence gives
\(\Diamond(b\AND P'_b(\NOT q))\IFF\Diamond(b\AND P_b(q))\). Again
\(\Diamond(b\AND P_b(q))\).

In both cases this contradicts \(\Box\NOT(b\AND P_b(q))\), which the
K-rule obtains from Lemma~\ref{lem:E} and the tautology
\(\NOT P_b(q)\IC\NOT(b\AND P_b(q))\). The remaining reasoning is
propositional, so the implication is a theorem of \HPE.
\end{proof}

\begin{lemma}[Necessary purity under the guard suffices]\label{lem:relative}
\(\Box\bigl(b\IC\kappa(b)\bigr)\IC F_b(Q_b(x))\).
\end{lemma}
\begin{proof}
Lemma~\ref{lem:guarded} gives the theorem
\((b\IC\kappa(b))\IC(b\IC C_b(Q_b(x)))\). The K-rule gives
\(\Box(b\IC\kappa(b))\IC\Box(b\IC C_b(Q_b(x)))\), and the consequent is
\(F_b(Q_b(x))\).
\end{proof}

Lettie's form of this lemma needs \(\Box\pi(P)\AND\Box\pi(P')\). With
the guard, \(\Box(b\IC\kappa(b))\) suffices, and that is something a base
theorem can provide.

\section{The twenty purity premises}\label{sec:premises}

We now state the finite assumptions. Besides the terms above, the final
argument uses two more closed terms that contain no constants:
\[
 Z:=\lambda a\lambda Y\lambda x.\bigl(a\AND Y(x)\bigr),\qquad
 \cmp:=\lambda Y\lambda S\lambda x.\,Y(S(x)).
\]
By \(\beta\)-conversion, \(Z\,a\,Y\,x\IFF(a\AND Y(x))\) and
\(\cmp\,Y\,S\,x\IFF Y(S(x))\). We use ten instances of Logical Purity,
\[
 \begin{aligned}
 &\mathrm{L1}:\Pure_t(\top),&&\mathrm{L2}:\Pure(\conj),&&
  \mathrm{L3}:\Pure(m),&&\mathrm{L4}:\Pure(\hW),&&
  \mathrm{L5}:\Pure(\hPhi),\\
 &\mathrm{L6}:\Pure(\hP),&&\mathrm{L7}:\Pure(\hPp),&&
  \mathrm{L8}:\Pure(\hQ),&&\mathrm{L9}:\Pure(Z),&&
  \mathrm{L10}:\Pure(\cmp),
 \end{aligned}
\]
and ten instances of Application,
\[
 \begin{aligned}
 &\mathrm{A1}:\App_{t,t\IC t},&&\mathrm{A2}:\App_{t,t},&&
  \mathrm{A3}:\App_{u,t},&&\mathrm{A4}:\App_{c,u\IC u},&&
  \mathrm{A5}:\App_{u,u},\\
 &\mathrm{A6}:\App_{c,u},&&\mathrm{A7}:\App_{c,t\IC u},&&
  \mathrm{A8}:\App_{t,u},&&\mathrm{A9}:\App_{t,u\IC u},&&
  \mathrm{A10}:\App_{u,u\IC u}.
 \end{aligned}
\]
Since \(u=t\IC t\), A1 and A8 are the same formula; it is listed twice
because it is used at two different points. Let \(\mathrm{LA}\) be the
conjunction of these premises. None of them is assumed necessary.

\begin{theorem}\label{thm:main}
\(\HPE\dentails\mathrm{LA}\IC\bigl(\PPu\IC\NOT\exists r\,\Uni(r)\bigr)\).
Hence, over C, \(\PPu\) with the Logical Purity and Application schemas,
all as ordinary premises, refutes the existence of a universal
proposition.
\end{theorem}

We prove the theorem in two stages: the purity premises produce the
guard \(a\) of \eqref{eq:goals}, and universality then yields a
contradiction relative to \(a\).

\section{The purity argument}\label{sec:purity}

From now on, the premises are \(\mathrm{LA}\) and \(\PPu\). Put
\[
 G:=\lambda p.\bigl(\Pure_t(p)\AND p\bigr),\qquad
 \Pres(\Theta):=\forall p\bigl((\Pure_t(p)\AND p)\IC\Theta(p)\bigr).
\]
Thus \(\Pres(\Theta)\) says that \(\Theta\) holds of every pure truth.
The predicate \(G\) contains the constant \(\Pure_t\). It is used only to
instantiate the quantifier in \(W\); its purity is never asserted.

\begin{lemma}[The guard is pure and true]\label{lem:closure-true}
Under L1--L4, A1--A5 and \(\PPu\): if \(\Theta\) is pure and
\(\Pres(\Theta)\), then \(a\) is pure and true.
\end{lemma}
\begin{proof}
Assume \(\Pure_u(\Theta)\) and \(\Pres(\Theta)\). We first show
\(\Cl(G)\), using \(G(p)\IFF(\Pure_t(p)\AND p)\).
\begin{itemize}
\item (Cl1): \(\top\) is true, and pure by L1.
\item (Cl2): if \(G(p)\) and \(G(q)\), then \(p\AND q\) is true. By L2,
  A1 and A2, \(\conj\,p\,q\) is pure, and it is \(p\AND q\) by
  \(\beta\)-conversion.
\item (Cl3): if \(G(p)\), then \(\Pres(\Theta)\) gives \(\Theta(p)\), and
  \(\Pure_u(\Theta)\), \(\Pure_t(p)\) and A2 make \(\Theta(p)\) pure.
\item (Cl4): if \(\pi(S)\), that is, \(\Pure_u(S)\), then \(\PPu\) and A3
  make \(\pi(S)\) pure, and \(\pi(S)\) is true.
\item (Cl5): if \(\pi(S)\) and \(\forall p(S(p)\IC G(p))\), then, since
  \(G(p)\IC p\), \(\forall p(S(p)\IC p)\), which is \(m(S)\). L3,
  \(\Pure_u(S)\) and A3 make \(m(S)\) pure.
\end{itemize}
Elimination for \(W\) at \(X:=G\) now gives \(W(p)\IC G(p)\), hence
\(W(p)\IC p\); generalizing \(p\), we obtain \(m(W)\). Next, \(\PPu\),
L4 and A4 make \(\hW\pi\) pure, and \(\Pure_u(\Theta)\) with A5 makes
\(W=\hW\pi\Theta\) pure, which is \(\pi(W)\). So
\(a=m(W)\AND\pi(W)\) is true.

Finally, L3, \(\pi(W)\) and A3 make \(m(W)\) pure; \(\PPu\), \(\pi(W)\)
and A3 make \(\pi(W)\) pure; and L2, A1 and A2 make their conjunction
\(a\) pure.
\end{proof}

\begin{lemma}[\(\Phi\) is pure and preserves pure truths]\label{lem:Phi}
Under the premises, \(\Pure_u(\Phi)\) and \(\Pres(\Phi)\).
\end{lemma}
\begin{proof}
\(\PPu\), L5 and A6 make \(\Phi=\hPhi\pi\) pure. Now let \(b\) be pure
and true. \(\PPu\), L6 and A7 make \(\hP\pi\) pure, and \(\Pure_t(b)\)
with A8 makes \(P_b=\hP\pi b\) pure. This is \(\pi(P_b)\). Likewise \(\PPu\), L7, A7, \(\Pure_t(b)\) and
A8 give \(\pi(P'_b)\). So \(\kappa(b)\) holds, and with \(b\),
\(\Phi(b)\).
\end{proof}

Now take \(\Theta:=\Phi\) in Section~\ref{sec:closure}. This fixes
\(W=\hW\pi\Phi\) and \(a=m(W)\AND\pi(W)\).

\begin{lemma}[Properties of the guard]\label{lem:guard}
\begin{enumerate}
\item[(a)] Under the premises, \(a\) is pure and true.
\item[(b)] \(\dentails\Box\bigl(a\IC\kappa(a)\bigr)\).
\item[(c)] \(\dentails F_a(Q_a(x))\) for every \(x\).
\end{enumerate}
\end{lemma}
\begin{proof}
Part (a) is Lemma~\ref{lem:closure-true} with Lemma~\ref{lem:Phi}. For
(b), \((a\IC\Phi(a))\IC(a\IC\kappa(a))\) is a theorem, by
\eqref{eq:kappa}; the K-rule and Lemma~\ref{lem:base} give
\(\Box(a\IC\kappa(a))\). For (c), apply Lemma~\ref{lem:relative} at
\(b:=a\) and use (b).
\end{proof}

The premises enter only in (a). Parts (b) and (c) are theorems of the
base logic, so no purity or truth claim has been necessitated. This is
\eqref{eq:goals}.

\section{The contradiction}\label{sec:final}

\begin{lemma}[Universality relative to the guard]\label{lem:step1}
Under the premises, if \(r\) is universal, then \(R_a(r,y)\) for every
\(y\).
\end{lemma}
\begin{proof}
Let \(Y\) be pure with \(Y(y)\), and put \(Z_Y:=Z\,a\,Y\). L9,
\(\Pure_t(a)\) and A9 make \(Z\,a\) pure, and \(\Pure_u(Y)\) with A5
makes \(Z_Y\) pure. By \(\beta\)-conversion, \(Z_Y(y)\) follows from
\(a\), which is true, and \(Y(y)\). Universality of \(r\) gives
\(R(r,y)\), and its instance at \(Z_Y\) gives \(\Diamond Z_Y(r)\).
Since \(\NOT(a\AND Y(r))\IC\NOT Z_Y(r)\) is a theorem, the K-rule turns
\(\Diamond Z_Y(r)\) into \(\Diamond(a\AND Y(r))\). Introduction for
\(R_a\) gives \(R_a(r,y)\).
\end{proof}

\begin{proof}[Proof of Theorem~\ref{thm:main}]
Assume the premises, and suppose that \(r\) is universal. Put
\(e:=Q_a(r)\). We first show \(R_a(e,\NOT e)\). Let \(Y\) be pure with
\(Y(\NOT e)\), and put \(Y':=\cmp\,Y\,Q_a\).

The operator \(Q_a=\hQ\pi a\) is pure: \(\PPu\), L8 and A7 make
\(\hQ\pi\) pure, and \(\Pure_t(a)\) with A8 makes \(Q_a\) pure. Hence
\(Y'\) is pure: L10, \(\Pure_u(Y)\) and A10 make \(\cmp\,Y\) pure, and
\(\Pure_u(Q_a)\) with A5 makes \(Y'\) pure.

Next, \(Y'(\NOT e)\). Lemma~\ref{lem:guard}(c) at \(x:=r\) makes \(e\)
\(a\)-firm. By symmetry \(\NOT e\) is \(a\)-firm too, so
Lemma~\ref{lem:B} gives \(Q_a(\NOT e)=\NOT e\). Leibniz's law turns
\(Y(\NOT e)\) into \(Y(Q_a(\NOT e))\), which is \(Y'(\NOT e)\).

Lemma~\ref{lem:step1} at \(y:=\NOT e\) gives \(R_a(r,\NOT e)\). Its
instance at \(Y'\) gives \(\Diamond(a\AND Y'(r))\). Since
\(Y'(r)\IFF Y(e)\), the K-rule gives \(\Diamond(a\AND Y(e))\).
Introduction for \(R_a\) now yields
\[
 R_a(e,\NOT e).
\]

But \(e\) is \(a\)-firm: \(\Box(a\IC C_a(e))\). By T and the truth of
\(a\), \(C_a(e)\) holds, and its first conjunct is
\(\NOT R_a(e,\NOT e)\). This contradiction shows that, under the
premises, no \(r\) is universal; existential Instantiation gives
\(\NOT\exists r\,\Uni(r)\). Separating \(\PPu\) from the other premises
gives the displayed theorem.

For the form over C, every theorem of \HPE\ is a theorem of C, and each
conjunct of \(\mathrm{LA}\) is an instance of Logical Purity or
Application. Modus ponens with the premises gives the conclusion.
\end{proof}

Where Lettie's proof uses Exhaustion to necessitate the true purity
claims \(\pi(P)\) and \(\pi(P')\), this proof necessitates only the base
theorem \(a\IC\Phi(a)\). The premises make the guard \(a\) pure and
true. Its purity lets the operators \(Z_Y\) and \(Q_a\) serve as pure
properties, and its truth discharges the guard after T.

\begin{corollary}
Let \(\mathrm{EF}:=\exists r\,\Fun_t(r)\), and let
\(\mathrm{H5}:=\Pure(\negop)\AND\App_{u,u}\AND\Rec\) be the hypotheses of
the Proposition of Section~1. Then
\[
 \HPE\dentails\mathrm{LA}\IC\bigl(\mathrm{EF}\IC(\mathrm{H5}\IC\NOT\PPu)\bigr).
\]
In particular, Logical Purity, Application, a unique fundamental
proposition, no fundamental entities at other types, and zeroary and
unary Recombination, as ordinary premises over C, refute \(\PPu\).
\end{corollary}
\begin{proof}
Under \(\mathrm{EF}\) and \(\mathrm{H5}\), the Proposition of
Section~1 gives a universal proposition, while \(\mathrm{LA}\) and
\(\PPu\) rule one out by Theorem~\ref{thm:main}. A unique fundamental
proposition gives \(\mathrm{EF}\). For the second claim, the premises
used are Logical Purity (L1--L10 and \(\Pure(\negop)\)), Application, a
unique fundamental proposition and \(\Rec\). The absence of fundamental
entities at other types and zeroary Recombination are not used. Adding \(\PPu\) to
these premises therefore derives a contradiction.
\end{proof}

\end{document}
